\textcolor{blue}{1. [25 min/30 puntos] Considera que un problema de manera secuencial se resuelve en un tiempo
de 2k nanosegundos, en donde k es el tamaño de la entrada (si el tamaño de la entrada es de
k = 100, entonces el cómputo toma 200 nanosegundos). Supongamos que este problema es $100\%$
paralelizable, si el número de hilos es de 8. ¿Cuál es el tamaño máximo de la entrada que se puede
procesar en 0.5 segundos ($5 \cdot 10^8$ nanosegundos)?} \textcolor{red}{Hint: utiliza nansosegundos en la ecuación.}\\

Lo primero que hacemos es convertir el tiempo en nanosegundos 
$$0.5 s = 5 \cdot10^8 ns$$

como el problema dice que es $100\%$ paralelizable usamos la Ley de Amdahl
$$S(n)=\frac{1}{(1-p)+\frac{p}{n}}$$

donde tenemos 
$$p=1$$
$$n=8$$

sustituimos

$$S(8)=\frac{1}{(1-1)+\frac{1}{8}} = 8$$

Tomando en cuenta el resultado anterior tomamos la fórmula 

$$S= \frac{T_s}{T_p}$$ 

y despejamos $T_p$, por lo que la fórmula nos queda como 

$$T_p = \frac{T_s}{S}$$

como $S=8$ tenemos y el tiempo secuencial es de $T_s= 2k$, tenemos $T_p=\frac{2k}{8}$

como el tiempo debe ser menos a $5\cdot10^8$ tenemos
$$\frac{2k}{8} = \frac{k}{4}$$
$$ \frac{k}{4} \leq 5\cdot10^8$$
$$k \leq 20 \cdot10^8= 2 \cdot10^9$$
$$= 2,000,000,000$$

En conclusión el tamaño máximo de la entrada sería de $2\cdot10^9$ \\